% !TeX root = ../papers/nc-p3.tex
\section{有限窓の反対座標積は十次元または十六次元}
\label{sec:finite-opposite-product-gap}
\begin{theorem}\label{thm:finite-opposite-product-gap}
四辺が有限な標準有限Frobenius窓で、反対二位置の座標を $U,W\subset D=H^0(\mathcal O_X)$ とする。
$U\otimes W\to D$ が単射でないこと、二移動での座標射が射影線型に同定されること、
その側の反対交換で四テンソルが不変になる表示を持つことは同値である。
この場合は $\dim UW=10$、それ以外では $\dim UW=16$ である。
点スキームの被約性、反対切断階数、印の存在、全次数の正則性を仮定しない。
\end{theorem}
\begin{proof}[階数一と二の双一次式]
定理\ref{thm:finite-square-chart}の有限完全交叉の証明により
$W^2=hU^2$ となる単元 $h$ があり、両座標の二次・三次の制限は単射である。
これらは有限Frobeniusの四移動定理により現在の候補にも適用できる。
$u b=0$ ($0\ne u\in U,0\ne b\in W$) なら $u(h^{-1}b^2)=0$ が三次多項式の恒等式となり矛盾する。
階数2の式を $u b=v c$ とする。各 $w\in W$ に対し $h^{-1}bw,h^{-1}cw$ を
$U$ 上の一意な二次式として持ち上げる。これらに $u,v$ を掛けたものは三次式として等しい。
一次式 $u,v$ は独立なので、ある $A(w)\in U$ により
\[
     h^{-1}bw=v A(w),\qquad h^{-1}cw=u A(w).
\]
$A$ は線型で、$bw=0$ は二次制限の単射性から $w=0$ を与えるため同型である。
隣の位置の座標と $W$ の間の任意の関係で $w$ を $A(w)$ に置き換えて $X$ 上で評価する。
得た $f$ は隣接座標積に属し、二等式から $uf=vf=0$ である。
その支持は $U$ 側の直線に入るので補題\ref{lem:finite-adjacent-word-support}により $f=0$。
グラフの一意性で $A$ が二移動を与え、有限辺の反転からテンソル対称性を得る。
逆に対称表示から二移動の同定が従う。実際二辺が反転し、有限グラフの両射影が同型である。
二移動の射影同定があれば、線束の自明化の単元を除いて $W=U$ であり積は十次元である。

\end{proof}
\begin{proof}[階数三は不可能]
$W$ の基底について積 $h^{-1}w_i w_j$ の二次持上げを対称行列 $P(x)$ に並べる。
その十成分は $U$ の全二次式の基底であり、$X$ 上では全二次小行列式がスキームとして零である。
混合双一次関係の係数を一次式の列 $m(x)$ とすれば、$P(x)m(x)=0$ が多項式として成り立つ。
実際各成分は元の式に $W$ の一文字を掛けた三次式である。逆も、$W$ に単元を取れることから従う。

階数3の関係があれば基底で $m(x)=(v_1,v_2,v_3,0)^t$ とできる。
$P$ の左上三行三列 $H$ は $H v=0$ により
\[
 H=[v]_\times^t C [v]_\times
\]
と一定の対称行列 $C$ で表せる。ここで $[v]_\times z=v\times z$。
この表示を確認する。第四変数で展開した各係数 $F(v,z)=z^t H_j(v)z$ は
$z\mapsto z+t v$ に不変で $v$ に関する次数が $j\le2$ である。
$F(v+t z,z)=t^{-2}F(v+t z,v)$ なので、$j<2$ の係数は零、$j=2$ では
左辺が $t$ に依存しない。よって $F$ は二つの剪断に不変で、独立な二列の対では
$v\times z$ だけに依存する（同じ外積を持つ基底は $\mathrm{SL}_2$ で移る）。
外積の非零座標ごとにその座標だけを分母に用いて二列を復元できる。
従ってこの関数は三空間の原点外で正則で、座標環の一意分解により原点にも延びる。
次数から外積の二次式となる。これで表示を得た。
$C$ が階数1以下なら $H$ は高々二独立な一次式の列の平方で、その六成分は独立になれない。
従って $q(v)=v^t\operatorname{adj}(C)v$ は非零である。
余因子行列の積の式から $\operatorname{adj}(H)=q(v)vv^t$ なので $X$ 上で $q(v)v_i v_j=0$。
三一次式の共通零点が $X$ にないなら直ちに $q=0$ となる。
あるならその点の極大イデアル $\mathfrak m$ をこの三式が生成し、$q$ はそこで
$\mathfrak m^2$ に消され、他では零である。$q\in U^2$ は隣接座標の平方全体と完全跡で直交する。
その平方は $\mathcal O_{X,p}/\mathfrak m^2$ に全射する（単元と、閉埋込みの一次座標差を用いる）。
ゆえに $q$ は零でなければならず、二次制限の単射性に反する。

\end{proof}
\begin{proof}[階数四の関係の二次行列]
階数4なら $m(x)=x$ として $P(x)x=0$ となる。一般階数3と仮定すると
$\operatorname{adj} P=\lambda(x)xx^t$ で、$\lambda$ は四次式である。
実際関数体での核が $kx$ なのでこの形になり、四変数に共通因子がないため係数は多項式である。
全三次小行列式とその各偏微分は $X$ 上で零である。偏微分は二次小行列式の多項式結合である。
$X$ の各点で単元な $x_i$ を選び $\lambda x_i^2$ とその偏微分を使うと、
全ての $\partial_j\lambda$ が $X$ 上で零となる。三次制限の単射性と標数零により $\lambda=0$ となり矛盾する。
従って $P$ の一般階数は2以下である。

一般階数2のとき $\wedge^2 P$ は対称な階数1の四次多項式行列で、$d(x)b(x)b(x)^t$ と書ける。
これは関数体で階数1の列を取り、多項式環で分母と全成分の共通因子を除き、
原始列への比例係数が多項式になることと対称性を使えばよい。
$b$ は一般点で $P$ の像の二平面を表す二形式で、$P x=0$ から $\iota_x b(x)=0$。
$\deg b=0$ は非零の定数二形式では不可能である。
$\deg b=1$ なら $b(x)=\iota_x\omega$ となる一定の三形式 $\omega$ がある。
実際偏極すると係数の $x$ 添字と縮約される添字も交代する。
四次元の三形式は非零の根基 $kz$ を持つので、一般の $x$ における $P$ の像二平面は
$z^\perp$ に入る。これは全二次式を基底とする $P$ に定数核を与えて矛盾する。
よって $\deg b=2$, $d$ は非零定数である。関数体上で非零な主二次小行列式は定数倍の平方なので、
標数零の代数閉基礎体より、階数2の二次式 $y^t P(x)y$ は関数体上で二一次式に分解する。
係数二次式は共通因子を持たないのでGaussの補題で多項式環上に分解が降りる。
一方の因子が $x$ に関して定数なら値域は固定一次式の倍だけになり、$P$ の全線型包に反する。
従って両因子は双一次で
\[
 y^t P(x)y=(y^t Sx)(y^t Tx).
\]
$P(x)\ne0$ は全射影点で成り立つので $S,T$ は可逆である。
一般の $Sx,Tx$ は独立で $P(x)x=0$ のため $S,T$ はともに交代行列である。
その鉛筆のPfaffianの根を取り $L=T-cS$ を階数2にする。両行列は比例せず、
$S$ は可逆なのでこのような有限な $c$ がある。
$L$ の像を根基に持つ階数2の交代行列 $J$ を取れば
$(Sx)^tJSx=(Tx)^tJSx=0$。従って $P(x)JSx=0$ は階数2の混合関係を与え、最初の証明が適用される。

一般階数1の場合も同じ原始列の議論で $P=d b b^t$ と書け、全成分の次数と独立性から $d$ は定数、
$b$ は可逆な一次列である。この場合 $X$ 上の $W$ 座標点は $[b(x)]$ に等しい。
実際二次Veroneseは閉埋込みで、$P$ は $X$ 上でも両者の同じ二次座標を与える。
これで全階数の場合が尽きる。
\end{proof}
